\section{Composición de conceptos}\label{sec:composicion}

\subsection{Una operación productiva y explícita}
Para componer conceptos uso la selección K20, que toma una variable \(x\), un dominio \(D\) y
una condición \(\varphi(x)\). Produce el conjunto
\[
  \{x\in D\mid\varphi(x)\}.
\]
El dominio puede ser un concepto indexado, un conjunto enumerado, una selección
anterior o un conjunto convencional declarado. Una lista no es un dominio de
selección en este fragmento; antes debe declararse el conjunto de sus miembros.
El cuerpo es una proposición. La variable queda vinculada solo dentro de él,
no dentro del dominio. Los parámetros restantes deben estar definidos fuera.

\KalaMostrar{\KalaSeleccion{x}{\KalaConcepto{3}}{\KalaPosee{\KalaVariable{x}}{2}}}
Lectura: sensores activos. C3 es el inventario de sensores y P2 fija el estado
operativo en un instante común. Se ha construido un concepto sin afirmar que
tenga miembros. Para afirmar que existe alguno se escribe:
\KalaMostrar{\KalaExiste{x}{\KalaConcepto{3}}{\KalaPosee{\KalaVariable{x}}{2}}}

\subsection{Patrones de formación}
\begin{description}[style=nextline]
\item[Base más propiedad.] Sensores activos: seleccionar sensores que satisfacen P2.
\item[Dos condiciones.] Sensores rojos y activos: unir dos predicados con K12 dentro
de la selección. T16 muestra la expresión completa.
\item[Exclusión dentro de un dominio.] Sensores que no son rojos: seleccionar
sensores con la negación de P1. Esto no crea automáticamente el concepto «azul».
\item[Restricción por relación.] Personas a las que les gusta el libro indicado:
seleccionar \(x\) en C1 con R7(\(x\), referencia del libro).
\item[Restricción cuantificada.] Personas a las que les gusta algún libro del
estudio: selección sobre C1, con un existencial sobre C2 en el cuerpo.
\end{description}

\KalaMostrar{\KalaSeleccion{x}{\KalaConcepto{1}}{\KalaExiste{y}{\KalaConcepto{2}}{\KalaGusta{\KalaVariable{x}}{\KalaVariable{y}}}}}

Otra disciplina usa la misma operación. Muestras líquidas bajo condiciones
experimentales declaradas:
\KalaMostrar{\KalaSeleccion{x}{\KalaConcepto{4}}{\KalaPosee{\KalaVariable{x}}{3}}}
C4 identifica las muestras y P3 incluye temperatura, presión y criterio de fase.

\subsection{Reglas de interpretación}
Se lee la base, se identifica la variable y se aplica la condición completa.
Dos signos pegados sin constructor no tienen una composición implícita. La
posición, el giro o un trazo añadido tampoco crean una operación no documentada.
El resultado puede ser vacío. Renombrar una variable vinculada conserva el
significado si se evitan capturas de otras variables.

\(\{x\in D\mid P(x)\land Q(x)\}\) y
\(\{z\in D\mid Q(z)\land P(z)\}\) tienen los mismos miembros en el mismo
contexto clásico. La definición del concepto y su nombre pueden conservar
información adicional; la coincidencia de miembros no autoriza borrar sus
identificadores de catálogo.

\subsection{Cómo abreviar sin perder significado}
Una definición K09 puede dar nombre a un término construido, además de las
definiciones matemáticas convencionales. Debe indicar que el resultado es un conjunto.
El nombre es una abreviatura local con expansión textual sin captura de variables.
No se admite autorreferencia ni ciclos de definiciones en este fragmento.
La introducción de un nuevo índice global exige registrar su expansión en el
catálogo. El texto M01 muestra una definición local y su uso posterior.
